\begin{helper}
Let \(H\) be a history cell, \(T\) a target event, \(U\) the terminal
coordinate set, and let \(o\) be a nodup terminal order with
\(\operatorname{toFinset}(o)=U\).  Suppose the terminal coordinates are
outside the recorded coordinates, and suppose both the recorded coordinates
and terminal coordinates lie in
\(\noderef{TwoBiteTerminalCoordinateUniverse}(m)\).  Thus all coordinates
used in the scan are tagged increasing non-loop base-edge coordinates, and
the recorded and terminal coordinate sets have compatible canonical
orientations.  Suppose \(T\subseteq H\), every
staged-open coordinate of a sample in \(H\) lies in \(U\), and every target
sample has open-pair count at least \(a\).  Assume the prefix-safety
hypothesis along \(o\): whenever two samples agree on the embedding, on the
recorded coordinates, and on all coordinates preceding \(e\) in \(o\), the
membership of \(e\) in staged-open pairs, touched pairs, same-color-closed
pairs, and pre-terminal recorded pairs is the same for the two samples.
Suppose also that the blue support label family \(\mathcal B\) consists of
\(u_B\)-element blue subsets of \(U\), and that every \(\omega\in T\) has a
label \(B\in\mathcal B\) equal to its present blue staged-open support.
Assume every \(\omega\in T\) has red present staged-open support of
cardinality \(u_R\).  Finally assume the red support decoder works uniformly
after a complete blue trace is fixed: whenever \(\omega\in T\) and
\(\beta_0\in H\) agree on all blue terminal coordinates of \(U\), there is a
red label \(J\) whose branch-decoded support from \(\beta_0\) is exactly the
present red staged-open support of \(\omega\).  Assume also that the edge
probability lies in the product-measure range \(0\le p\le1/2\).

Then there is a finite type \(\mathcal T\) of branch/blue-trace labels, a
finite set \(\mathcal P\) of transcripts, a blue-trace map
\(\operatorname{blueTrace}:\mathcal T\to\{\hbox{subsets of }U\}\), a
representative branch \(b(\beta)\in H\) for every label
\(\beta\in\mathcal T\), a beta-indexed red decoder
\[
\operatorname{redSupport}(\beta,J)
  =\operatorname{branchRedSupport}(b(\beta),J),
\]
and cell events
\[
G_{B,J,\beta,\tau}\subseteq \{\hbox{samples}\},
\]
indexed by \(B\in\mathcal B\), \(J\in\mathcal J\),
\(\beta\in\mathcal T\), and \(\tau\in\mathcal P\), such that:

For every \(\beta\in\mathcal T\), the set
\(\operatorname{blueTrace}(\beta)\) is a subset of \(U\), consists only of
blue coordinates, and records the actual terminal blue truth table of the
representative branch: for every blue coordinate \(e\in U\),
\[
X_e(b(\beta))=1
  \quad\Longleftrightarrow\quad
e\in\operatorname{blueTrace}(\beta).
\]
The label is unique for this terminal blue truth table: if
\(\beta,\beta'\in\mathcal T\) have the same membership on every blue
coordinate of \(U\), then \(\beta=\beta'\).  Thus extra nonterminal or red
coordinates cannot distinguish two labels with the same complete terminal
blue trace.

First, the cells cover \(T\cap H\): every \(\omega\in T\cap H\) lies in some
\(G_{B,J,\beta,\tau}\) with \(B\in\mathcal B\),
\(\beta\in\mathcal T\), and \(\tau\in\mathcal P\).

Second, whenever \(B\in\mathcal B\), \(J\in\mathcal J\),
\(\beta\in\mathcal T\), \(\tau=(P_\tau,A_\tau,Z_\tau)\in\mathcal P\), and a
sample \(\omega\) lies in \(G_{B,J,\beta,\tau}\), the following properties
hold:
\[
b(\beta)\in H,\qquad
\operatorname{redSupport}(\beta,J)
  =\operatorname{branchRedSupport}(b(\beta),J),
\]
and \(b(\beta)\) has the complete terminal blue trace
\(\operatorname{blueTrace}(\beta)\); that is, for every blue coordinate
\(e\in U\),
\[
X_e(b(\beta))=1
  \quad\Longleftrightarrow\quad
e\in\operatorname{blueTrace}(\beta).
\]
\[
|\operatorname{redSupport}(\beta,J)|=u_R,\qquad
\operatorname{redSupport}(\beta,J)\subseteq U,
\]
and this decoded support consists only of red coordinates;
\[
P_\tau,A_\tau\subseteq U,\qquad P_\tau\cap A_\tau=\emptyset,
\]
\[
B\cup\operatorname{redSupport}(\beta,J)\subseteq P_\tau,\qquad
Z_\tau\subseteq A_\tau,
\]
and \(P_\tau,A_\tau\) encode the complete blue terminal trace:
for every blue coordinate \(e\in U\),
\[
e\in\operatorname{blueTrace}(\beta)\quad\Longleftrightarrow\quad
e\in P_\tau,
\qquad
e\notin\operatorname{blueTrace}(\beta)\quad\Longleftrightarrow\quad
e\in A_\tau .
\]
\[
\max(0,a-u_R-u_B)\le |Z_\tau|.
\]
Moreover, relative to any sample \(\omega\) lying in the cell, the label
\(\beta\) records the complete blue terminal trace of \(\omega\): for every
blue coordinate \(e\in U\),
\[
X_e(\omega)=1
  \quad\Longleftrightarrow\quad
e\in\operatorname{blueTrace}(\beta),
\]
and the transcript forces every coordinate of \(P_\tau\) present and every
coordinate of \(A_\tau\) absent in \(\omega\).

The construction also exposes the combinatorial support data needed by the
later fixed-\((B,J)\) product-tree summation.  There is a predicate saying
which formal pairs \((B,J,\beta,\tau)\) are realized, and an
assignment-compatibility predicate for terminal assignments.  The realized
predicate is not merely a zero/nonzero tag: if
\(\mathsf{Realized}(B,J,\beta,\tau)\) holds, then the realized pair satisfies
the red-support cardinality and color facts, the transcript subset and
disjointness facts, \(B\cup\operatorname{redSupport}(\beta,J)\subseteq
P_\tau\), \(Z_\tau\subseteq A_\tau\), the complete blue-trace equivalences,
and the selected-absence lower bound displayed above.  Compatible
assignments contain \(P_\tau\), avoid \(A_\tau\setminus Z_\tau\), and agree
with the complete blue trace on all unselected blue coordinates.  Functional
scan disjointness says that one terminal assignment cannot be compatible
with two distinct realized \((\beta,\tau)\) for the same \(B,J\).

The node does not assert fixed-\((B,J)\) residual scan completeness,
residual owner uniqueness, residual leaf ownership, or residual mass
normalization.  Those fixed-\((B,J)\) reveal-order and product-tree facts are
supplied downstream, at the fixed-\((B,J)\) residual support layer where the
scan, released blocks, and leaf certificates are explicit.  Non-realized
formal pairs have no cell samples.

Thus this node remains the branch/transcript interface: it constructs the
branch labels, transcripts, realized geometry, complete-assignment
compatibility, functional disjointness, and non-realized-cell exclusion.  It
deliberately leaves the residual scan and released-block partition to the
later fixed-\((B,J)\) support nodes.
\end{helper}
\begin{proof}
For a target sample \(\omega\), choose its blue support label
\(B\in\mathcal B\).  Let \(\operatorname{Blue}(\omega)\) be the complete set
of blue terminal coordinates of \(U\) which are present in \(\omega\).  Among
the finitely many complete blue traces which occur for target samples, choose
one representative branch \(b(\beta)\in H\) for each trace.  The label
\(\beta\) is this complete blue trace together with the chosen representative
for that trace, and \(\operatorname{blueTrace}(\beta)\) is exactly the
terminal blue coordinate set of that trace.  Hence
\(\operatorname{blueTrace}(\beta)\subseteq U\), every member of
\(\operatorname{blueTrace}(\beta)\) is blue, and membership in
\(\operatorname{blueTrace}(\beta)\) agrees with the value of
\(b(\beta)\) on every terminal blue coordinate.  Because only one
representative label is chosen for each complete terminal blue truth table,
two labels with the same membership on all blue coordinates of \(U\) are
equal.  The branch partition in
\(\noderef{FixedSetHistoryCellAdaptiveBranchPartition}\) supplies the cover
and disjointness of the complete blue trace classes.
These representative branches are chosen as global returned data: for every
label \(\beta\), \(b(\beta)\in H\), independently of whether a later cell
with label \(\beta\) is nonempty.

For the target sample \(\omega\), take the unique label \(\beta\) whose
blue trace is \(\operatorname{Blue}(\omega)\).  By construction,
\(\omega\) and \(b(\beta)\) agree on every blue terminal coordinate.  The
uniform red support-label hypothesis applied to this representative branch
therefore gives a red label \(J\) such that
\[
\operatorname{branchRedSupport}(b(\beta),J)
\]
is exactly the red present staged-open support of \(\omega\).  Define
\(\operatorname{redSupport}(\beta,J)\) to be this branch-decoded support.
This definition is global in the two labels: for every \(\beta\) and every
red label \(J\),
\[
\operatorname{redSupport}(\beta,J)
  =\operatorname{branchRedSupport}(b(\beta),J).
\]
Thus the red package is still allowed to depend on the representative branch
chosen for the complete blue trace, but there is only one branch label for
that trace, so later summations do not acquire multiplicity from several
representatives with the same blue values.

For a nonempty cell, the cell event remembers the target sample, the chosen
representative branch, and the assigned red label.  The assigned-label
equality says that the decoded support is the red present staged-open support
of that sample.  Therefore the decoded support has cardinality \(u_R\), is
contained in \(U\), consists only of red coordinates, and every coordinate in
it is both staged-open and present.  The cardinality assertion uses the
explicit fixed red-count hypothesis on target samples.  The containment and
red-color assertions follow from the definition of red present staged-open
support and from staged-open containment in \(U\).  The blue label hypothesis
gives the same facts for \(B\), with cardinality \(u_B\) and blue color.

Run the deterministic prefix scan along the order \(o\).  The scan records a
triple \(\tau=(P_\tau,A_\tau,Z_\tau)\): coordinates forced present,
coordinates forced absent, and the selected absent coordinates.  In the
returned transcript, put the present supports
\(B\cup\operatorname{redSupport}(\beta,J)\) into \(P_\tau\), put the whole
complete blue terminal trace \(\operatorname{blueTrace}(\beta)\) into
\(P_\tau\), and put every blue terminal coordinate outside that trace into
\(A_\tau\).  The selected absences \(Z_\tau\) are also inserted into
\(A_\tau\).  Since \(B\) is contained in the blue trace of the sample, the
red support is red-only, and \(Z_\tau\) is disjoint from
\(B\cup\operatorname{redSupport}(\beta,J)\), the two forced sets are
disjoint.  Hence for every nonempty cell witnessed by \(\omega\) one has
\(P_\tau,A_\tau\subseteq U\), \(P_\tau\cap A_\tau=\emptyset\),
\[
B\cup\operatorname{redSupport}(\beta,J)\subseteq P_\tau,\qquad
Z_\tau\subseteq A_\tau.
\]
Moreover the construction gives, for every blue \(e\in U\),
\[
e\in\operatorname{blueTrace}(\beta)\Longleftrightarrow e\in P_\tau,
\qquad
e\notin\operatorname{blueTrace}(\beta)\Longleftrightarrow e\in A_\tau .
\]

The canonical absence selection lemma
\(\noderef{FixedSetHistoryCellCanonicalAbsenceSelection}\) applies to the red
and blue present supports.  Its prefix-safety input is exactly the
prefix-safety hypothesis of this node: for a split \(o=\mathrm{pre}\, e\,
\mathrm{suffix}\), any comparison sample agreeing on the embedding, recorded
coordinates, and all coordinates in \(\mathrm{pre}\) has the same
staged-open, touched, same-color-closed, and pre-terminal-recorded status at
\(e\).  The remaining inputs are the nodup order, the equality
\(\operatorname{toFinset}(o)=U\), the containment of all staged-open
coordinates in \(U\), and the facts just proved that the red and blue support
sets are exactly the present staged-open coordinates already accounted for.
The target lower bound
\[
a\le |\operatorname{StagedOpen}(\omega,\varepsilon,I)|
\]
and the fixed support sizes \(u_R,u_B\) give
\[
\max(0,a-u_R-u_B)\le |Z_\tau|.
\]
By construction of the complete blue trace, for every blue terminal
coordinate \(e\in U\), \(X_e(\omega)=1\) exactly when
\(e\in\operatorname{blueTrace}(\beta)\).  By
construction of the transcript, all coordinates in \(P_\tau\) are present and
all coordinates in \(A_\tau\) are absent in any sample lying in the cell.
The sample's own labels and transcript place it in one of the displayed
cells, giving the cover of \(T\cap H\).

Finally record the product-tree support data.  For fixed \(B,J,\beta,\tau\),
declare
\[
\mathsf{Realized}(B,J,\beta,\tau)
\]
to mean exactly that the cell \(G_{B,J,\beta,\tau}\) is nonempty.  Thus, if
the formal pair is not realized and a sample \(\omega\) lay in the cell, then
\(\omega\) itself would witness nonemptiness, a contradiction.  This proves
the required zero-contribution statement for non-realized formal pairs before
any summation over formal indices is made.
Conversely, if \(\mathsf{Realized}(B,J,\beta,\tau)\) holds, choose a witness
\(\omega\in G_{B,J,\beta,\tau}\).  Applying the nonempty-cell geometry
already proved above to this witness gives the red-support size and color,
the containments in \(U\), \(P_\tau\cap A_\tau=\emptyset\),
\(B\cup\operatorname{redSupport}(\beta,J)\subseteq P_\tau\),
\(Z_\tau\subseteq A_\tau\), the complete blue-trace equivalences, and
\(\max(0,a-u_R-u_B)\le |Z_\tau|\).  This is the exported
realized-geometry certificate used by the mixed normalization package; it is
not inferred from the one-way non-realized exclusion.

For a terminal assignment \(A\subseteq U\), define
\(\operatorname{Compat}(A;B,J,\beta,\tau)\) as a stronger
internal certificate than the three forward consequences exported in the
statement.  First,
\[
P_\tau\subseteq A,\qquad A\cap A_\tau=\emptyset .
\]
The displayed exported avoidance
\[
A\cap(A_\tau\setminus Z_\tau)=\emptyset
\]
then follows immediately, since \(A_\tau\setminus Z_\tau\subseteq A_\tau\).
This distinction is essential: disjointness from
\(A_\tau\setminus Z_\tau\) alone would not say anything about selected
coordinates in \(Z_\tau\).  Second, for every blue terminal coordinate
\(e\in U\setminus Z_\tau\),
\[
e\in A\quad\Longleftrightarrow\quad
e\in\operatorname{blueTrace}(\beta).
\]
Third, the complete blue trace, including selected coordinates, is the
branch trace:
\[
\forall e\in U,\quad e\hbox{ blue}\Longrightarrow
  \bigl(e\in A\Longleftrightarrow
    e\in\operatorname{blueTrace}(\beta)\bigr).
\]
Fourth, the realized compatibility certificate is single-owner: if another
realized \((\beta',\tau')\) has \(P_{\tau'}\subseteq A\), the full
disjointness \(A\cap A_{\tau'}=\emptyset\), and the same complete blue
terminal trace on \(A\), then \(\tau'=\tau\).  Thus compatibility is not
merely a span condition; it contains the full selected-coordinate absence
needed before applying the realized single-owner comparison.  This
realized-only comparison is built into the local compatibility certificate
so that no total branch assignment has to be chosen.

The exported forward consequences are now checked one by one.  The first
component gives \(P_\tau\subseteq A\).  Full disjointness from \(A_\tau\)
gives the stated weaker disjointness from \(A_\tau\setminus Z_\tau\).  The
blue-trace components give the complete blue trace, and hence also the
unselected blue-trace agreement stated in the helper.  If one terminal
assignment \(A\) is compatible with two realized cells \((\beta,\tau)\) and
\((\beta',\tau')\), the two complete-blue-trace clauses and
branch-functionality give \(\beta=\beta'\).  The compatibility certificate
for \((\beta,\tau)\) can then be applied to the second realized cell because
the second compatibility certificate supplies exactly the missing premises:
\[
P_{\tau'}\subseteq A,\qquad A\cap A_{\tau'}=\emptyset,
\]
together with the same complete blue trace on \(A\).  Therefore
\(\tau'=\tau\).  Hence a complete terminal assignment is compatible with at
most one realized \((\beta,\tau)\) for the fixed \(B,J\).

The fixed-\((B,J)\) released residual blocks and residual leaf masses are
intentionally not constructed in this node.  Those data belong to the
downstream fixed-count summation layer, which combines this node's realized
geometry, compatibility, and functional disjointness with the separate
fixed-\((B,J)\) residual scan and leaf support.

The non-realized-cell exclusion is immediate from the definition of
\(\mathsf{Realized}\).  If a formal pair is not realized and some sample
\(\omega\) lay in \(G_{B,J,\beta,\tau}\), then \(\omega\) would witness that
the cell is nonempty, contradicting non-realization.  This completes the
branch/transcript interface consumed by the later fixed-count summation.
\end{proof}
